\section{Frozen factorial trading experiment}
\label{sec:factorial}
\paragraph{Models and paired conditions.}
We evaluate Qwen2.5-Coder 7B and 14B at the immutable revisions recorded in the protocol,
with sampling seeds 11, 23, and 37. For each model, seed, and basket, the agent receives
base access, knowledge access (K), numerical-tool access (T), or both (K+T). Each decision
allows at most eight turns and 1,024 generated tokens per turn at temperature 0.1.
Condition order is seeded. Each cell has 44 scheduled decisions, giving
$2\times2\times3\times4\times44=\FTDecisions$ decisions. No success or profit threshold
selects the analysis population.

\paragraph{Assets, data, and timing.}
The original basket contains SPY, QQQ, EFA, EEM, TLT, IEF, GLD, DBC, NVDA, AVGO, ORLY,
and FAST. These assets and their period were already used in development. The transfer
basket was declared before downloading its outcomes: XLB, XLC, XLE, XLF, XLI, XLK, XLP,
XLRE, XLU, XLV, XLY, and SHY. It changes the instruments while retaining the calendar
and substantial exposure overlap. Both use the January 3, 2025--September 25, 2026
scoring window, with 434 sessions and available warmup history.

All cells receive the same total-return price index constructed from the frozen quotes
and corporate actions, restricted to the decision date. The prompt states that the
prices are already adjusted. Targets execute after the next session's close-to-close
return; rebalancing opportunities occur every ten sessions. Missing or invalid submissions
retain drifted holdings. Portfolios are long-only, cash earns zero, and proportional
cost is 5 bps per traded side. The ledger includes the first trade's fee and initial
capital. Fixed strategy references use the same timing and cost conventions.

\paragraph{Availability and accounting checks.}
Before full inference, CPU qualification checked access restrictions, future-read denial,
adjusted-price parity, algorithm execution, and valid submissions. A single development
date per model and condition tested transport; a non-submission was retained rather
than converted into a quality gate. Protocols, sources, data hashes, and the analysis
plan were frozen before the full run. We subsequently verified all decision receipts
and independently recomputed every cell's net-asset-value path with a dollar-holdings
ledger. This verification adds no model decisions or selection rules.

\paragraph{Endpoints.}
The predeclared primary endpoint is cumulative net return from initial capital. The agent
prompt asks for net risk-adjusted return; Sharpe and drawdown are secondary reported
measures rather than substitutes chosen after observing results. For a model, basket,
and seed, let $B,K,T,F$ denote cumulative returns in the four cells. Paired effects are
$F-B$, the knowledge marginal effect $[(K-B)+(F-T)]/2$, the tools marginal effect
$[(T-B)+(F-K)]/2$, and interaction $F-K-T+B$, in percentage points. We average these
within model and basket over the three declared seeds.

Risk, exposure, turnover, submissions, successful document reads, numerical calls,
exact output adoption, failed turns, tokens, and generation time are secondary outcomes.
We also reprice the same fixed decisions at 0, 10, and 25 bps without simulating a
behavioral response. Circular calendar-block resampling uses lengths 10, 20, and 40,
with 2,000 draws each; indices are shared across all cells and seeds within each group.
The resulting 95\% intervals describe dependence sensitivity on this path.

\begin{table*}[t]
\centering\small
\setlength{\tabcolsep}{4pt}
\begin{tabular}{llrrrrrrrr}
\toprule
Basket & Model & Base & K & T & K+T & Full$-$Base & K effect & T effect & Interaction \\
\midrule
Original & 7B & 28.10 & 34.73 & 54.56 & 33.80 & 5.70 & -7.07 & 12.76 & -27.39 \\
Original & 14B & 40.65 & 44.91 & 28.95 & 39.92 & -0.73 & 7.62 & -8.35 & 6.71 \\
Transfer & 7B & 10.73 & 15.17 & 13.20 & 12.42 & 1.69 & 1.83 & -0.14 & -5.21 \\
Transfer & 14B & 9.67 & 10.21 & 4.70 & 14.37 & 4.70 & 5.11 & -0.41 & 9.13 \\
\bottomrule
\end{tabular}
\caption{Mean cumulative net return (percent) and paired effects (percentage points), over three sampling seeds on the same historical calendar. All planned decisions, including missing submissions, contribute. K denotes skill-document access; T denotes numerical-tool access with calling contracts. Original assets were used in development; transfer changes assets within the same period.}
\label{tab:factorial-main}
\end{table*}


\section{Trading results and observed use}
\paragraph{Mixed return differences.}
Table~\ref{tab:factorial-main} reports every group mean. Relative to base, full access
changes original-basket returns by \FTOriginalSevenDifference{} points for 7B and
\FTOriginalFourteenDifference{} for 14B. The transfer differences are
\FTTransferSevenDifference{} and \FTTransferFourteenDifference{} points. Every full-access
group falls below the equal-weight rebalancing reference
(Table~\ref{tab:factorial-baselines}). Both positive and negative seed-level differences
occur within every model--basket group (Table~\ref{tab:factorial-seeds}). All reported
calendar-block intervals span zero (Table~\ref{tab:factorial-intervals}); no general
performance advantage is established.

\paragraph{Access rarely becomes use.}
Across the complete batch, \FTSkillReads{} successful skill-document reads and
\FTAlgorithmDecisions{} decisions with successful numerical calls were observed
(Table~\ref{tab:factorial-use}). In particular, algorithm use concentrates in the transfer
7B full-access cell; no tools-only cell makes a successful numerical call. Exact matching
between submitted weights and a tool output is reported separately from calling it.
The zero-read result makes interpreting K as consumed financial knowledge untenable.
The return contrasts estimate offered access, including changes to prompts and behavior.

\paragraph{Failures, exposure, and costs.}
Submission denominators include all scheduled decisions, and failed turns remain in the
record. A missed rebalance can leave a profitable or losing position in place, so returns
do not directly measure interface competence. Tables~\ref{tab:factorial-risk}
and~\ref{tab:factorial-cost} report exposure, drawdown, turnover, token use, generation time,
and transaction-cost sensitivity. Appendix~\ref{app:factorial} and the aggregate evidence
retain all seed-level results; none are selected for successful tool use or positive return.
